%% The following is a directive for TeXShop to indicate the main file
%%!TEX root = diss.tex

%=======================================
\chapter{Higher-dimensional Lattice Walk Models}
\label{ch:HigherD}
%=======================================

We begin the exploration of higher-dimensional models by making some useful definitions. Here are the main ingredients for a model set in $\ze^d$.

\begin{itemize}
	\item A fixed finite set $\mc S$ of steps for travel between points in the space. It may involve arbitrarily-many finitely-sized steps with multiplicities at least one. Its corresponding Laurent polynomial is denoted by $S$.
	\item A region $\mc D \subseteq \ze^d$, where the walks are free to grow but not venture outside of the boundaries. This may be the whole space, an orthant, a semi-infinite strip, finite boxes, and so on.
	\item A subregion $\mc D_0 $ of $\mc D$, where interactions take place. To each point $p$ there will be a corresponding label $a_p$ whose power we can increment upon visits to $p$. Common subregions $\mc D_0$ include the origin, the axes in two dimensions, various higher-dimensional hyperplanes, a cone, an infinite strip, the sides of a box and so on.
	\item A set $\mc Q$ of initial walks that we can evolve. It has a corresponding generating polynomial $q$.
\end{itemize}

We will be growing lattice walks from an initial condition $\mc Q$, with steps $\mc S$, within $\mc D$ in $\ze^d$, while counting interactions within $\mc D_0$. The four sets determine the model completely.

\begin{defn} A lattice walk model $\mc M$ is characterized by the following four parts: $\mc S$, $\mc D$, $\mc D_0$, $\mc Q$.
\end{defn}

\begin{defn} Define the \classic $d$-dimensional lattice walk model to be the model with small steps in the positive orthant in $\ze^d$. The interacting \classic model is the \classic $d$-dimensional model with interactions with the hyperplanes $x_k = 0$ bounding the non-negative orthant.
\end{defn}

\begin{exa} Consider one of the one-dimensional models from the previous chapter. This is an example of a an interacting \classic one-dimensional model. Here $\mc D$ is the set of non-negative integers, $\mc D_0$ consists of the point at the origin, where the marking variable is $a$, $\mc Q$ consists of the empty walk. Finally, $\mc S$ is one of the subsets of $\{\um 1, 0, 1\}$.
\end{exa}

\begin{exa} Consider classic models with interactions in two dimensions. Here $\mc D = \{(x_1,x_2)\in \ze^2\mid x_1,\, x_2 \ge 0\}$ and $\mc D_0$ consists of marking points on the non-negative axes --- $x_1$ with label $a_2$, since $x_2 = 0$ here, and $x_2$ with label $a_1$, since $x_1 = 0$. Note that the origin is marked with both labels. As in the previous example, $\mc Q$ consists of the empty walk and $\mc S$ is some set of small steps, where each step has coordinates of size at most 1.
\end{exa}

We also need to keep track of all the properties of the lattice walks in order to enumerate them, so we will use a generating function that tracks quantities like length and ending position. Since the sets $\mc S$ and $\mc Q$ are finite, the number of walks of any given length will be finite as well. Therefore, the coefficients of $Q$ will be Laurent polynomials.
\begin{defn} Let $\mc M$ be a lattice walk model and $Q$ its associated generating function counting lattice walks' lengths with $t$, ending locations with $\row xd$, and interactions with $\mc D_0 $ with some multi-index variable $a$. Then we write
\begin{align}
	Q = Q(\row xd,a;t) = \sum_{k=0}^{\infty} q_k(\row xd, a)t^k\,,
\end{align}
where the coefficients $q_k$ are Laurent polynomials in the $x$ and $a$ variables. 
\end{defn}

\begin{exa} In this example we will examine the simplest kind of model --- one where $\mc D$ is the whole space, there are no interactions and $\mc Q$ consists of the empty walk. Since the lattice walks may pick any step from the step set $\mc S$, the walks of length one are enumerated via $t S$, the walks of length two, via $t^2S^2$, and so on, where $S = S(\row xd)$. Then the full generating function is
\begin{align}
	Q(\row xd;t) = 1 + St + S^2t^2 + \cdots = \frac{1}{1-St}\,.
\end{align}
\end{exa}

\begin{exa} Consider the two-dimensional \classic model with $\mc S = \{$NE, NW, SE, SW$\}$, that is, $S = (\ol x_1 + x_1)(\ol x_2 + x_2)$. The only walk of length one is enumerated by $x_1x_2t$. When we grow the walk, it may take any of the four available steps and we have
\begin{align}
	Q = 1 + x_1x_2t + (a_1 + x_1)(a_2 + x_2)t^2 + \cdots\,.
\end{align}
We may extract coefficients too. For example, the walks ending on the $x_1$ axis are given by setting $x_2 = 0$ or extracting the $x_2^0$ term. That is,
\begin{align}
	[x_2^0]Q = 1 + a_2(a_1 + x_1)t^2 + \cdots\,.
\end{align}
The walks returning to the origin are given by
\begin{align}
	[x_1^0x_2^0]Q = 1 + a_1a_2t^2 + \cdots\,.
\end{align}
\end{exa}

While it is possible to handle the analysis of higher-dimensional \classic models without extra notation, it will greatly help us to have it. In fact, notation conventions paved the way to finding the functional equations for higher-dimensional models and later, the operators. With this thesis, we aim to make this kind of notation more common in the enumeration of lattice walks. Typically, in the multivariate setting we see multi-index notation where a bolded or vectorial form denotes the collection of variables involved, but such notation is not enough when working with restricted geometries for regions and marking subregions. We have the following notation. 

\begin{defn}Let $\oo := (0,\ldots,0)$ and $\pm \one := (\pm 1,\ldots, \pm 1)$, where we have as many terms in these sequences as needed in the specific context. Also, let $[n] = \{1,\ldots, n\}$.
\end{defn}

\begin{defn} Let $\row xd$ be some indeterminates and let $I = (\row im)$ and $J = (\row jm)$ be sequences of positive integers. Define the monomial $x_I^J$ as
\begin{align}
	x_I^J = x_{i_1}^{j_1}\cdots x_{i_m}^{j_m}\,.
\end{align}
Then let the coefficient and term extraction operators be defined as
\begin{align}
	[x_I^J] F= [x_{i_1}^{j_1}\cdots x_{i_m}^{j_m}]F = F_I^J \qandq [[x_I^J]] = x_I^J [x_I^J]\,.
\end{align}
For the monomial, we have the special cases $x_I^{\oo} = 1$ and $x_I^{\one} = x_I$ and $x_I^{\um\one} = 1/x_I = \ol x_I$. For the coefficient extraction, we have the special cases $F_I^J = F$ if $I$ is empty, and $F_I^J = F_J$ if $I$ is the full set of variables.
\end{defn}

\begin{exa}In the previous example we had $[x_2^0]Q = Q_2^0$ and $[x_1^0x_2^0]Q = \Qf{0,0}$. Furthermore, we have $\Sf{\pm 1,\pm1}  = 1$ and 
\begin{align}
	S_1^{\um 1} &= [\ol x_1](x_1 + \ol x_1)(x_2 + \ol x_2) = x_2 + \ol x_2\,.
\end{align}
\end{exa}

Because there are different variables involved, there may be some ambiguity in the result, depending on the order of extraction. We will always perform coefficient extractions with respect to the $x$ or $a$ variables by first expanding any series in $t$, and then applying the extraction on the coefficients. The following notation makes this precise.
\begin{defn} Suppose that $\Gamma$ is an operator involving $x$ and $a$ variables but not $t$. We apply $\Gamma$ on power series $F = \sum_{k=0}^\infty f_k t^k$ in the following way:
\begin{align}
	\Gamma F = \sum_{k=0}^\infty \Gamma (f_k) t^k\,.
\end{align}
\end{defn}

\begin{exa} In the past we have performed inclusion-exclusion to arrive at functional equations for the two-dimensional problems. This is the usual procedure carried out in problems so far, like in \cite{bousquet2010walks}, but when it comes to generalizing to models with interactions, the process is difficult and prone to error. For a view of difficulties involved in not using the operator notation, see \cite{tabbara2014exact}, where the authors had to wrestle with case-by-case analysis. Instead, we use operators to derive the functional equations cleanly. As in the one-dimensional case with generating function $Q = \sum_{k=0}^{\infty} q_k t^k$, we may grow the set of walks by performing the multiplication $Sq_k$. Then we assume that an operator $\Gamma$ exists such that it cleans up the terms in $Sq_k$. We obtain the functional equation $(1-t\Gamma S)Q = 1$ again and all that is left to be determined is the operator $\Gamma$ itself. We can make and observations that will allow us to find it easily. The one-dimensional operator $1- [[\ol x]]$ is idempotent, that is $E^2 = E$, and combinatorially-speaking, once we remove certain terms,  we do not remove them again. In two dimensions this means that with the use of one-dimensional operators, we may remove terms with negative $x_1$ coordinate first, and ones with negative $x_2$ coordinate second, without performing inclusion-exclusion. That is, the operator $E$ here is given by $E_1E_2$ where $E_k = 1 - [[\ol x_k]]$. Multiplying these out, we find that 
\begin{align}
	E = 1 - [[\ol x_1]] - [[\ol x_2] + [[\ol x_1 \ol x_2]]
\end{align}
and the operator iteslf encodes the inclusion-exclusion. The situation is analogous for the marking operator $L = L_1 L_2$, where $L_k = 1+b_k[x_k^0]$. Again, the operator $L$ is invertible and analogously, we may write the functional equation as $(L^{\um 1} - tES)Q = \ol a_1 \ol a_2$. If we perform the coefficient extractions, we will find the equation
\begin{align}
	Q = 1 + t SQ - t\ol x_1S_1^{\um 1}Q_1^0 - t\ol x_2S_2^{\um 2}Q_2^0 + t\ol x_1\ol x_2\Sf{\um 1,\um 1}\Qf{0,0}\,.
\end{align}
\end{exa}

Contrast this notation to the one in \cite[p.~10]{bousquet2010walks}. The problem quickly runs into notational difficulties for $d$-dimensional problems, but that can be dealt with by writing the equation as a sum over subsets of $\{1,\ldots, d\}$, like it is done for the non-interacting case in Equation 5.1 in \cite[p.~58]{melczer2014variants}. This method runs into trouble when we start considering interactions. We take the functional equation for the non-interacting case and modify it to add interactions. This analysis has been done before in Section 3 of \cite[p.~7]{tabbara2014exact} and Section 3.2 of \cite[p.~11]{wong2014confining}, where the authors add in terms like $(a_1-1)T$, where $T$ is some expression involving extractions of $Q$ and $S$, but we will not go into the messy details. Instead, we will use the operator notation to handle the problem in a more efficient manner.
%=======================================
\section{Using Operators to Derive the Functional Equation}
%=======================================

As we have seen in the previous example, it is not difficult to generalize the derivation of the functional equations if we use operators. The derivation of general case for $d$ dimensions is analogous to the one in the example, so we will omit it most of it. In $d$ dimensions, the interactions take place on the hyperplanes where $x_k = 0$, so we will use the labels $a_k$ to mark those interactions. The operators here are $L = \mon Ld$ and $E = \mon Ed$, where order does not matter. We need to compute the quantity $L^{\um 1}(1)$, so we have
\begin{align}
	L^{\um 1}(1) = \sum_{I\subseteq [d]}(-1)^{|I|}c_I = \ol a_{[d]}\,.
\end{align}
The higher-dimensional functional equation for the interacting \classic model is given as follows.
\begin{prop} The \classic interacting $d$-dimensional lattice walk model has a full generating function $Q$ satisfying the functional equations
\begin{align}
	(1-tLES)Q &= 1\,,\\
	(L^{\um 1}-t ES)Q &= \ol a_{[d]}\,,
\end{align}
where $L$, $E$, $L^{\um 1}$ are of the form $\Gamma = \mon{\Gamma}{d}$.
\end{prop}
Previously, we have had to cope with coefficient extractions and back-substitution to find the expanded form of the second functional equation \cite[p.~9]{tabbara2014exact} and \cite[p.~12]{wong2014confining}). We can skip this process by using operators. In general, the second functional equation has many less terms. To see this, we take a look at the operator $L_kE_k$. We have
\begin{align}
	L_kE_k = 1 + b_k[x_k^0]-\ext{\ol x_k}
\end{align}
and if we apply the operator to an equation, we increases the number of terms in the equation by a factor of 3. We have $3^d$ terms relating to $SQ$ by the end of applying $LE$. On the other hand, once we apply the operator $L^{\um 1}$, we introduce $2^d$ terms relating to $Q$, but since we only have $E$ applied on $SQ$, we would only have at most $2^d$ terms. We have exactly $2^d$ terms $Q_{I}^\oo$, one for each subset $I \subseteq [d]$. In this case, we can write out the functional equation explicitly. First,
\begin{subequations}
	\label{eq:ExpExt}
	\begin{align}
		L^{\um 1} Q &= (1-c_1[x_1^0])\cdots(1-c_d[x_d^0])Q = \sum_{I\subseteq [d]} (-1)^{|I|} c_I Q_I^\oo\,,\\
		 ESQ &= \sum_{I\subseteq [d]} (-1)^{|I|} \ol x_I [\ol x_I](SQ)= \sum_{I\subseteq [d]} (-1)^{|I|} \ol x_I S_I^{\um \one}Q_I^\oo\,,
	\end{align}
\end{subequations}
where the last equality follows from the fact that there are no negative powers of $x_k$ in $Q$ for any $k$. Then the functional equation is
\begin{align}
	\sum_{I \subseteq [d]}(-1)^{|I|}(c_I-t \ol x_IS_{I}^{\um \one})Q_I^\oo = \ol a_{[d]}\,.
\end{align}
We let $K_I = c_I-t \ol x_IS_{I}^{\um \one}$, where, if $I =\emptyset$, we have the empty products and empty extractions, leaving us with $K_{\emptyset} = K$. Finally, we can write the master equation in $d$ dimensions.
\begin{thm} For subsets $I$ of $\{1,\ldots,d\}$, let $K_I = c_I-t \ol x_I S_{I}^{\um \one}$. Then the \classic interacting $d$-dimensional lattice walk model has a generating function satisfying 
\begin{align}
	\label{eqn:Master}
	\sum_{I\subseteq [d]}(-1)^{|I|} K_I Q_I^\oo= \ol a_{[d]}\,.
\end{align}
\end{thm}
\noindent Equation 5.1 of \cite[p.~58]{melczer2014variants} is the non-interacting version of our equation, and it can be written as
\begin{align}
	KQ - \exx\sum_{\substack{\emptyset \neq I \subseteq [d]}}\exx(-1)^{|I|}\ol x_IS_I^{\um \one}Q_I^{\oo} = 1\,.
\end{align}

%=======================================
\section{Coefficient Extractions}
%=======================================

In this section we will see how the extraction and marking operators interact with basic coefficient extraction to produce quicker coefficient extraction of the functional equation. That is, it is faster to use the operator form of the functional equation in order to extract coefficients than it is to use the explicit form. However, this section is not essential for our understanding of the enumeration of lattice walk models. The main role of the following analysis is to provide a shortcut for manual calculations. After all, we need to obtain functional equations for $Q_I^\oo$ for $I \subseteq [d]$ like in the one-dimensional case.
%=======================================
\subsection{Using Operator Notation For Fast Coefficient Extraction}
%=======================================
 We will find the $x_I^\oo$ term of the functional equation to obtain further equations for the functions $Q_I^\oo$. The operator $[x_i^k]$ commutes with operators $M_j = L^{\um 1}$ and $E_j$ when $j\neq i$, so we only need to consider the action on $M_i$ and $E_i$. Similarly to \autoref{eq:OpExt1} and \autoref{eq:OpExt2}, in the two-dimensional setting we have
\begin{align}
	[x_i^k]M_i = \ol a_i^{\delta(k,0)}[x_i^k] \qandq [x_i^k]E_i = [x_i^k]\,.
\end{align}
Therefore, for $k\ge 1$ the operator $[x_i^k]$ eliminates both $M_i$ and $E_i$ from the picture, while for $k = 0$ it removes $E_i$ but leaves a residual $\ol a_i$ behind for $M_i$. Therefore, if we extract $[x_I^\oo]$ from the functional equation for $I \subseteq [d]$, we have
\begin{align}
	\begin{split}
		[x_I^\oo]a_{[d]} =\ol a_{[d]} &= [x_I^\oo](M_{[d]} - t E_{[d]}S)Q \\
		&= \left(\ol a_{I}M_{[d]\backslash I}[x_I^\oo] - t E_{[d]\backslash I}[x_I^\oo]S\right)Q\\
		&= [x_I^\oo](\ol a_{I}M_{[d]\backslash I}- t E_{[d]\backslash I}S)Q\,,
	\end{split}
	\shortintertext{so that}
	\label{eq:OpExt3}
	\ol a_{[d]\backslash I} &= [x_I^\oo](M_{[d]\backslash I}- a_{I}t E_{[d]\backslash I}S)Q\,.
\end{align}
When we extract a general term, we use the extraction $[x_I^\oo x_J^P]$, where $P$ is a nonempty set of integers greater than zero and $I \cup J = K$ with no overlaps. Through the same process as before, we find the equation
\begin{align}
	[x_I^\oo x_J^P]( M_{[d]\backslash K} - a_I t E_{[d]\backslash K} S)Q = 0\,.
\end{align}

%=======================================
\subsection{Further Functional Equations}
%=======================================
If we want to find equations similar to the master functional equation, we will need to perform the extraction explicitly like in \autoref{eq:ExpExt}. First,
\begin{align}
	[x_I^\oo] SQ  = S_I^\oo Q_I^\oo + R_I\,,
\end{align}
where $R_I$ is a collection of $S_I^{\um P} Q_I^{P}$ terms that involve sets $P$ with some nonzero entries. We can use this term inside the functional equation as
\begin{align}
	\begin{split}
		\ol a_{[d]\backslash I} &= M_{[d]\backslash I}[x_I^\oo]Q- a_{I}t E_{[d]\backslash I}[x_I^\oo](SQ)\\
		& = M_{[d]\backslash I}Q_I^\oo- a_{I}t E_{[d]\backslash I}(S_I^\oo Q_I^\oo + R_I)\\
		& = (M_{[d]\backslash I} - a_{I}t E_{[d]\backslash I}S_I^\oo)Q_I^\oo  -a_{I}t E_{[d]\backslash I}R_I
	\end{split}
\end{align}
so we can rewrite it as
\begin{align}
	(M_{[d]\backslash I} - a_{I}t E_{[d]\backslash I}S_I^\oo)Q_I^\oo = a_I(\ol a_{[d]} + t R_I')\,,
\end{align}
where $R_I' =  E_{[d]\backslash I}R_I$. We compute the following terms
\begin{subequations}
	\begin{align}
		M_{[d]\backslash I}Q_I^\oo &= \sum_{J \subseteq [d]\backslash I} (-1)^{|J|}c_J Q_{I,J}^{\oo,\oo}\,,\\
		a_I t E_{[d]\backslash I}S_I^\oo Q_I^\oo &= \sum_{J \subseteq [d]\backslash I}  (-1)^{|J|}a_I t\ol x_J S_{I,J}^{\oo,\um \one}Q_{I,J}^{\oo,\oo}\,,
	\end{align}
\end{subequations}
and the functional equation can be written as
\begin{align}
	\sum_{J \subseteq [d]\backslash I}  (-1)^{|J|}(c_J - a_I t\ol x_J S_{I,J}^{\oo,\um \one})Q_{I,J}^{\oo,\oo} = a_I(\ol a_{[d]} + t R_I')\,.
\end{align}
The form of this expression suggests the following notation.
\begin{defn} For integer sets $I$ and $J$ that do not share any elements, we write
\begin{align}
	K_J(I) = c_J - a_I t\ol x_J S_{I,J}^{\oo,\um\one}\,,
\end{align}
where if $I = \emptyset$, then $K_J(I) = K_J$.
\end{defn}
With this notation, we can write the $2^d$ equations for $Q_I^\oo$ and $I\subseteq [d]$ in simplified form.
\begin{lem} Let $I$ be a subset of $\{1,\ldots,d\}$. The generating function $Q_I^{\oo}$ for the $d$-dimensional interacting \classic lattice walk model satisfies the functional equation
\begin{align}
	\label{eqn:FullEqns}
	\sum_{J \subseteq [d]\backslash I}  (-1)^{|J|}K_J (I) Q_{I,J}^{\oo,\oo} = \ol a_{[d]\backslash I} + a_I t  R_I'\,,
\end{align}
where the $R_I'$ term involves only higher-order coefficient extractions from $Q$.
\end{lem}
This is a clean way to write down all the equations for $Q_I^\oo$ in a unifying way. To compute the $R_I'$ terms, we first determine what $R_I$ is. It is a sum of terms $S_I^{\um P} Q_I^{P}$ for various non-negative $P$ with size $|I|$ whose elements are not all 0. Then we compute $R_I' = E_{[d]\backslash I}R_I$. A key observation is that the terms $Q_I^P$ are power series in the remaining variables, so they do not have non-negative coefficients. Let $sq = S_I^{\um P} Q_I^{P}$. Then
\begin{align}
	E_j sq &= sq - \ol x_j(sq)_j^{\um 1} = sq - \ol x_js_j^{\um 1}q_j^0\
	\shortintertext{and thus,}
	E_{K} sq &= \sum_{J\subseteq K} (-1)^{|J|}\ol x_J s_J^{\um \one} q_J^\oo\,.
\end{align}
This allows us to derive the following explicit form of $R_I'$.
\begin{prop} The functions $R_I'$ are given by
\begin{align}
	R_I' =  \sum_{J \subseteq [d]\backslash I}(-1)^{|J|}\sum_{P > \oo} S_{I,J}^{\um P,\um \one}Q_{I,J}^{P,\oo}\,.
\end{align}
\end{prop}
Of course, when $d = 1$ or $d = 2$, we do not need the full generality of the formula, but for higher dimensions, it would be quite difficult to write the equations by hand without it.

%======================================
\subsection{Kernel Forms}
%======================================

We can write out the functional equations in a way that highlights the kernel $K$. For example,
\begin{align}
	(1-tLES)Q =(1-t(LE-1 + 1)S)Q = KQ + t(1-LE)SQ\,,
\end{align}
where $1-LE$ is an operator consisting of coefficient extractions only. Here is another way we could have expanded the operator form.
\begin{align}
	(1-tLES)Q = LEKQ + (1-LE)Q =LEKQ + (1-L)Q\,,
\end{align}
where $EQ = Q$ since $Q$ has no terms that need removing. Something interesting happens when we set all the $a_k$'s to 1 in the last equation. This is equivalent to removing the marking operator. The equation becomes $ E(KQ) = 1$ and this form suggests that we use the following substitution
\begin{align}
	Q = \frac{W}{\mon xd K}\,.
\end{align}
Then we multiply both sides of the functional equation by $\mon xd$ and obtain
\begin{align}
	\mon xd = \mon xd E(KQ) = \mon xd E(\ol x_1 \cdots \ol x_d W) = E' W\,,
\end{align}
where  $E' = E_1'\cdots E_d'$ and 
\begin{align}
	E_k' = x_kE_k\ol x_k = x_k(1-\ext{\ol x_k})\ol x_k = 1 - [x_k]^0\,.
\end{align}
\begin{prop} Let $W = \mon xd KQ$, where $Q$ is the generating function for walks without interactions in the positive orthant. Then
\begin{align}
	E' W = \sum_{I\subseteq [d]}(-1)^{|I|} W_I^\oo= \mon xd \,.
\end{align}
\end{prop}
For example, in two dimensions we have
\begin{align}
	W(x_1,x_2) - W_1^0(x_2) -W_2^0(x_1) + W_{1,2}^{0,0} = x_1x_2\,.
\end{align}
Thus, we can write the solution to the problem as
\begin{align}
	Q = \frac{1}{\mon xd K}\left(\mon xd-\sum_{\emptyset \neq I\subseteq [d]}(-1)^{|I|} W_I^\oo\right)\,.
\end{align}
In two dimensions we have
\begin{align}
	Q(x_1,x_2) = \frac{x_1x_2 + W_1^0(x_2) + W_2^0(x_1)-W_{1,2}^{0,0}}{x_1x_2(1-t S)}\,.
\end{align}

\begin{exa} Consider the one-dimensional model on the non-negative line with $S = \ol x + x$. Here the solution for the interacting version of the model is given by
\begin{align}
	Q = \frac{t(\ol \sigma - \ol x)}{(1-t S)(a\ol \sigma t- b)},
\end{align}
so when we set $a = 1$ and simplify the expression, we find that the generating function for the non-interacting model is given by
\begin{align}
	\wh Q = \frac{x - \sigma}{x(1-t S) }\,,
\end{align}
as predicted by the analysis of this section. Here $W_0 = \sigma$.
\end{exa}

Of course, each $W_I^\oo$ depends on $Q_I^\oo$, so we could have seen this relationship starting with the master functional equation for $Q$ itself.
Interestingly enough, the calculation points to a differential equation. Clearly, $W$ satisfies the differential equation $\del_1\del_2 W = 1$ and a similar relationship exists for the $d$-dimensional models.
\begin{prop} The generating function for a $d$-dimensional \classic model, satisfies the PDE
\begin{align}
	\mon \del d (\mon xd KQ) = 1\,.
\end{align}
\end{prop}
The result is not difficult to generalize. If we have walks with larger steps, we will have a similar partial differential equation where the orders of variables and orders of the derivatives will be higher. If we have interactions, we will also have higher order terms.
%=======================================
\section{First Solutions}
%=======================================
%=======================================
\subsection{Formal Neumann Series}
%=======================================
So far we have obtained the equation $(1-t\Gamma S)Q = 1$, where $\Gamma = LE$. Formally, we can expand this operator as Neumann series and express the generating functions as
\begin{align}
	Q = \sum_{k=0}^{\infty}t^k (\Gamma S)^k(1)\,,
\end{align}
where $(\Gamma S)^k(1)$ indicates repeated application of $\Gamma S$. This allows us to compute $Q$ iteratively, where the next $t$-coefficient is the previous one with $\Gamma S$ applied to it. This is a formalization of the naive brute force approach of growing lattice walks from smaller ones: by taking a set of walks of length $n$, extending them by 1 step, removing the ones that end outside of the quarter plane, and finally by marking the ones that interact with the hyperplanes $x_i = 0$ at their endpoint. We have not been able to find closed forms for the repeated operator applications in general. Here is an example that illustrates the use of the operators.
\begin{exa} Let $S = x_1 + \ol x_1 + x_2 + \ol x_2 = S_1(x_1) + S_2(x_2)$. We start with $(\Gamma S)^0(1) = 1$. Then
\begin{align}
	(\Gamma S)(1) &= LE(S_1 + S_2) =L(x_1 + x_2) = a_2x_1 + a_1x_2\,,\\
	\begin{split}
		(\Gamma  S)^2(1) &= LE(S_1 + S_2)(a_2x_1 + a_1x_2) \\
		&= LE(a_2S_1x_1 + a_1x_2S_2 + a_2x_1S_2 + a_1x_2S_1)\\
		&=a_1a_2^2 + a_2^2x_1^2 + a_1a_2^2 + a_1^2x_2^2 + (a_1 + a_2)x_1x_2\,.
	\end{split}
\end{align}
\end{exa}

%=======================================
\subsection{Hadamard Products}
%=======================================

We take a moment to see a result that falls out of the operator notation. In Section 5 of \cite[p.~14]{bostan2014three}, the authors consider the special form of the step polynomial $S(x_1,x_2) = S_0(x_1) + S_1(x_1)S_2(x_2)$ for the non-interacting case. Here we will find a similar result for the simpler case $S = S_1(x_1)S_2(x_2)$ but with interactions, and its generalizations to arbitrarily-many dimensions.
\begin{exa} When $S = S_1(x_1)S_2(x_2)$ the operator $\Gamma S$ factors as
\begin{align}
	 \Gamma S = \Gamma_1\Gamma_2S_1S_2 = (\Gamma_1S_1)(\Gamma_2S_2)\,.
\end{align}
This is because $\Gamma_2$ does not interact with $S_1$ and the operators commute. Then
\begin{align}
	 (\Gamma S)^k = (\Gamma_1\Gamma_2S_1S_2)^k = (\Gamma_1S_1)^k(\Gamma_2S_2)^k
\end{align}
and we use the Neumann series to find
\begin{align}
	Q = \sum_{k = 0}^{\infty} t^k (\Gamma S)^k = \sum_{k = 0}^{\infty} t^k (\Gamma_1S_1)^k(\Gamma_2S_2)^k = \left(\sum_{k = 0}^{\infty} t^k (\Gamma_1S_1)^k\right) \odot  \left(\sum_{k = 0}^{\infty} t^k (\Gamma_2S_2)^k\right)\,,
\end{align}
where $\odot$ denotes the Hadamard product of two power series with respect to variable $t$. By using
\begin{align}
	Q_i(a_i,x_i;t)= \sum_{k = 0}^{\infty} t^k (\Gamma_i S_i)^k
\end{align}
and reversing the iteration, we reach the functional equations $(1-t\Gamma_i S_i)Q_i = 1$ for each $i$. The functions $Q_i$ are solutions to one-dimensional problems with step sets $S_i$ and $Q$ is expressed as
\begin{align}
	Q(a_1,a_2,x_1,x_2;t) = Q_1(a_1,x_1;t) \odot Q_2(a_2,x_2;t)
\end{align}
\end{exa}

\begin{exa} Once we know the form of the solution, we can prove that it works by a method that does not use the series solution. Suppose that $Q_i$ is the solution to the functional equation
\begin{align}
	(L_i^{\um 1} - t E_i S_i)Q_i = \ol a_i
\end{align}
for $i = 1$ and $i = 2$. Then we claim that $Q = Q_1 \odot Q_2$ is a solution to the functional equation
\begin{gather}
	\begin{align}
		(L^{\um 1} - t ES)Q = \ol a_1\ol a_2\,.
	\end{align}
	\shortintertext{To this end, we compute}
	\begin{align}
		\begin{split}
			(L^{\um 1}-t ES)Q &= L^{\um 1} (Q_1\odot Q_2) - t ES(Q_1\odot Q_2)\\
			& = (L^{\um 1}_1Q_1\odot L^{\um 1}_2Q_2) - t((E_1S_1Q_1)\odot (E_2S_2Q_2))\,.
		\end{split}
	\end{align}
\end{gather}
Then we substitute $ L^{\um 1}_i Q_i = \ol a_i + t E_i S_i Q_i$ into the first term and expand it to find
\begin{align}
	\begin{split}
		 (L^{\um 1}_1Q_1\odot L^{\um 1}_2Q_2) &= (\ol a_1 + t E_1 S_1 Q_1)\odot(\ol a_2 + t E_2 S_2 Q_2)\\
		 & = \ol a_1\odot\ol a_2 + \ol a_1\odot (t E_2 S_2 Q_2) + \ol a_2\odot(t E_1 S_1 Q_1) \\
		 &\quad + (t E_1 S_1 Q_1)\odot(t E_2 S_2 Q_2))\,.
	\end{split}
\end{align}
Since we are working with power series,
\begin{align}
	\ol a_i\odot (t E_j S_j Q_j) &= \ol a_i [t^0]t E_j S_j Q_j = 0
\end{align}
and we are done.
\end{exa}

\begin{lem} Let $Q$ be the generating function for an interacting \classic model, with $S = S_1(x_1)S_2(x_2)$, and $Q_i$, with step polynomial $S_i$. Then
\begin{align}
	Q(x_1,x_2,a_1,a_2;t) = Q_1(x_1,a_1;t)\odot Q_2(x_2,a_2;t)\,.
\end{align}
\end{lem}
This identity will serve as the base case for the following general case.
\begin{thm} Let $Q$ be the generating function for an interacting \classic model with $S = \mon Sr$, where $S_i$ and $S_j$ do not share variables when $i\neq j$, and $Q_i$, for a model with step polynomial $S_i$. Then
\begin{align}
	Q = Q_1 \odot \cdots \odot Q_r\,.
\end{align}
\end{thm}

\begin{exa} Consider the two-dimensional model with diagonal steps, that is, $S = (x_1 + \ol x_1)(x_2 + \ol x_2)$. We will look at this case in more detail in \autoref{ch:OKM}. The one-dimensional problems are the same and have solutions
\begin{gather}
	\begin{align}
		\widetilde{Q}_0 = \frac{\sigma}{at-b\sigma}\qandq \widetilde{Q} = \frac{t(x-\sigma)}{x(1-t(x + \ol x)(at-b\sigma)}\,,
	\end{align}
	\shortintertext{where}
	\begin{align}
		\sigma = \frac{1-\sqrt{1-4t^2}}{2t}\,.
	\end{align}
\end{gather}
We can expand and simplify $\widetilde{Q}_0$ and $\widetilde{Q}$ to find
\begin{align}
	\widetilde{Q}_0 &= \frac{2-a-a\sqrt{1-4t^2}}{2(1-a + a^2t^2)}\,,\\
	\widetilde{Q} &= \frac{((a-2)t + (1-a)x + 2 a x t^2) + (at + (1-a)x)\sqrt{1-4t^2}}{2(1-a + a^2t^2)(x - t - t x^2)}\,.
\end{align}
We have the expansion
\begin{gather}
	\begin{align}
		\widetilde{Q} = 1 + x t + (a + x^2)t^2 + x(1 + a + x^2)t^3 + \cdots
	\end{align}
	\shortintertext{and finally,}
	\begin{align}
		\begin{split}
			Q(x_1,x_2,a_1,a_2;t) &= \widetilde{Q}(x_1,a_1;t) \odot \widetilde{Q}(x_1,a_2;t)\\
			&=1 + x_1x_2t + (a_1 + x_1^2)(a_2 + x_2^2)t^2\\
			&\quad + x_1x_2(1 + a_1 + x_1^2)(1 + a_2 + x_2^2)t^3 + \cdots\,.
		\end{split}
	\end{align}
\end{gather}
However, obtaining an explicit function via the Hadamard product seems to be rather messy. The function is not algebraic since the unmarked variant is not algebraic \FIXME{cite something}. However, the function $Q$ is D-finite since both functions in the product are.
\end{exa}

%========================================
\subsection{Reducible Models and Groups For Walks With Small Steps}
%========================================
Here is another observation that will allow us to find solutions to models in an easier way. As observed in \FIXME{cite walks with small steps, page so and so}, if the model has some steps that cannot be used in the construction of lattice walks, then we can safely remove them from the step set, and we would still have the same generating function enumerating the lattice walks. This procedure works for any step set and any regional constraints. 

\begin{defn}For a model with a step set $\mc S$, let the minimal step set $\mc S'$ for be the set of steps that can eventually be used by lattice walks. A model is reducible if $\mc S'$ is a proper subset of $S$.
\end{defn}

\begin{lem} Whether we use $\mc S$ or $\mc S'$, the full generating function for the model is the same.
\end{lem}

Next, it is easier to find the generating function for a model if we have more functional equations because this may help us cancel out some of the terms. One way to find more equations is to apply a transformation on the tuple $(\row xd)$ that leaves intact as many parts of the functional equation as possible but still adds some information in the form of a non-equivalent equation. We could map each $\row xd$ to some functions separately and then examine the group generated by such maps. We often eliminate the kernel in front of $Q$ in order to find equations for the extracted forms of $Q$, so it is useful to consider transformations that preserve the kernel. This is a fruitful approach in the non-interacting case, and for its application on two-dimensional models, see Section 4.2 \cite[p.~11]{bousquet2010walks}. We will call this group $G(K)$, although it is not the full group of rational transformations preserving the kernel. For example, a lot of the time $(x_1,x_2) \mapsto (x_2,x_1)$ is not in the group, but it is a valid transformation preserving the kernel. If a transformation on $(\row xd)$ preserves the kernel $K = 1 -t S$, then it must also preserve the step polynomial $S$. We may write $S$ as
\begin{align}
	S = S_k^{\um 1} \ol x_k + S_k^0 + S_k^1 x_k
\end{align}
for some $k$ from 1 to $d$. We send $x_k \mapsto y_k$ and solve $S(x_k) = S(y_k)$ for $y_k$ to determine the possible transformations we can perform. After some simplification, we obtain the equation
\begin{align}
	(S_k^{\um 1} - S_k^1 x_k y_k)(y_k-x_k) = 0\,.
\end{align}
We have three cases depending on which ones of $S_k^{\um 1}$ and $S_k^1$ are zero.
\begin{itemize}
	\item[1.] If both are zero, the equation is trivially-true and all transformations preserve the kernel.
	\item[2.] If only one of them is zero, then $y_k = x_k$, so only the identity transformation preserves the kernel.
	\item[3.] If neither one is zero, we have $y_k = x_k$ or $y_k$ is the involution $y_k = \ol x_k S_k^{\um 1}/S_k^1$.
\end{itemize}

\noindent Case one and two are mostly trivial and we find properties of the generating function right away.
\begin{itemize}
	\item[1.] If $S_k^{\um 1} = 0$ and $S_k^1 = 0$, then $S = S_k^0$ and the model takes all of its steps inside the $x_k = 0$ hyperplane. Therefore, the model is equivalent to a lower-dimensional model with no $x_k$ coordinate. Suppose that the generating function in the latter case is $q(t)$. Then the generating function in the former case is given by $Q(x_k, a_k;t) = q(a_kt)$, where the rest of the variables have been suppressed from the notation. We can also find this solution from the functional equation directly. Let $ M = M' M_k$ and $E = E' E_k$. Then by \autoref{eq:OpExt3} and the observation that $[x_k^0]SQ = SQ$ and $[x_k^0]Q = Q$, we have
	\begin{align}
		(M' - a_kt E' S)Q &= \ol a_{[d]\backslash \{k\}}\,,
	\end{align}
	This is just an equation for a $(d-1)$-dimensional model with an extra scaling of $t$ by $a_k$.
	\item[2.] If $S_k^{\um 1} \neq 0$ and $S_k^1 = 0$, then the model is reducible to one with $S = S_k^0$, that is, the model reduces to case 1.
	\item[3.] If $S_k^{\um 1} = 0$, $S_k^0 = 0$ and $S_k^1 \neq 0$, then the walks in model leave the $x_k = 0$ hyperplane and never return. Here we have an essentially lower-dimensional problem as well, since we have $S = S_k^1 x_k$. If the generating function for the lower-dimensional problem is $q(t)$, then the generating function for this model is $Q(x_k, a_k;t) = q(x_k t)$, where the rest of the variables have been suppressed from the notation. Similarly to case 1, we may derive the result from the functional equation directly.
\end{itemize}
All the other cases are more complicated and require more careful analysis. We can define a group $G_S$ of the nontrivial transformations.
\begin{defn} For a specific step polynomial $S$, let $G = G_S = G(K)$ be the group generated by the nontrivial transformations $g_k$.
\end{defn}

It is easy to see that there are infinitely many such groups, just consider $S = \prod_{k=1}^d (x_k + \ol x_k)$ with $g_k(x_k) = \ol x_k$. The group here has order $2^d$. The orders of the finite groups set in 2 dimensions are 1, 2, 4, 6 and 8, from Theorem 3 in \cite[p.~8]{bousquet2010walks}. If we were to consider walks with multiple copies of some steps, we may obtain even larger group orders as seen in Section 7 of \cite[p.~11]{kauers2014walks}. The $d$-dimensional groups defined above are Coxeter groups --- a generalization of the dihedral groups encountered in dimension 2 lattice walk models. Coxeter groups are defined as follows.
\begin{defn} The group
\begin{align}
	G = \langle \row rd \mid(r_ir_j)^{m_{i,j}} = 1\rangle
\end{align}
is known as a Coxeter group. Here $m_{i,i} = 1$ and $m_{i,j} \ge 2$ for all $i\neq j$. Also $m_{i,j}$ could be infinite.
\end{defn}
We call these $r$'s reflections and when $m_{i,j} = 2$, they commute. Also, $m_{i,j} = m_{j,i}$. The \emph{rank} of $G$ is the number of generators in this specific representation of $G$. Often $G$ has more than one representation with different numbers of generators. We can draw a graph for each Coxeter group by the following procedure.
\begin{itemize}
	\item Label vertices with the generators of $G$.
	\item Connect the vertices $g_i$ and $g_j$ if $m_{i,j} \ge 3$.
	\item Label the edges with $m_{i,j}$ if $m_{i,j}\ge 4$.
\end{itemize}
If the graph has different components, the Coxeter group factors as a direct product of Coxeter groups, represented by the individual components. We can see this phenomenon in terms of the step polynomial $S$. 
\begin{prop} If the step polynomial $S$ factors into two disjoint step polynomials $T$ and $R$ as
\begin{gather}
	\begin{align}
		S = T(\row xk)&R(x_{k + 1},\ldots,x_d)\,,
	\shortintertext{then the group $G_S$ factors as}
		G_S = G_T& \times G_R\,,
	\end{align}
\end{gather}
where $G_S$, $G_T$ and $G_R$ are the groups for the step polynomials $S$, $T$ and $R$, respectively.
\end{prop}
\begin{proof} For a fixed integer $1 \leq i \leq k$,
\begin{gather}
	\begin{align}
		S_i^j = [x_i^j]S = R[x_i^j]T = R T_i^j
	\end{align}
	\shortintertext{and therefore,}
	\begin{align}
		g_i(x_i) = \ol x_i\frac{T_i^{\um 1}R}{T_i^1R} = \ol x_i\frac{T_i^{\um 1}}{T_i^1}\,.
	\end{align}
\end{gather}
Similarly, if $k + 1 \leq j \leq d$, then $g_j(x_j) = \ol x_j R_i^{\um 1}/R_i^1$. Furthermore, $g_i$ and $g_j$ commute since they do not affect each other's variables. Since $i$ and $j$ are arbitrary, the generators from different sets commute and the Coxeter graph of the group factors in two, that is, the group itself factors.
\end{proof}


%=======================================
\section{Specializing to Two Dimensions}
%=======================================
Until the end of the chapter we will focus on the two-dimensional models with interactions. We can set $d=2$ in \autoref{eqn:Master} and \autoref{eqn:FullEqns} to find the two-dimensional master functional equations.
\begin{cor} The generating function for a two-dimensional \classic interacting model satisfies the following functional equations
\begin{align}
	\begin{split}
	KQ - K_1 Q_1^0 - K_2 Q_2^0 + K_{1,2} \Qf{0,0} &= \ol a_1\ol a_2\,,\\
		(1-ta_1 S_1^0)Q_1^0 - (c_2-ta_1\ol x_2S_{1,2}^{0,\um 1})Q_{1,2}^{0,0} - a_1 t(S_1^{\um 1}Q_1^1-\ol x_2S_{1,2}^{\um 1,\um 1}Q_{1,2}^{1,0}) &=\ol a_2\,,\\
		(1-ta_2 S_2^0)Q_2^0 - (c_1-ta_2\ol x_1S_{1,2}^{\um 1,0})Q_{1,2}^{0,0} -a_2t(S_2^{\um 1}Q_2^1-\ol x_1S_{1,2}^{\um 1,\um 1}Q_{1,2}^{0,1}) &=\ol a_1\,,\\
		(1-a_1 a_2 t \Sf{0,0})Q_{1,2}^{0,0} - ta_1a_2t(\Sf{\um 1,\um 1}\Qf{1,1} + \Sf{\um 1,0}\Qf{1,0} + \Sf{0,\um 1}\Qf{0,1})&= 1\,.
	\end{split}
\end{align}
\end{cor}

%=======================================
\subsection{First Examples in the Plane}
%=======================================

The \classic two-dimensional lattice walk models have been investigated thoroughly but the interacting version has not. We want to extend the \classic model methods to the interacting setting or create new ones. Since the generating functions for the interacting models need to reduce to those for the non-interacting models, we do not expect to find solutions for every model. In fact, some of the \classic models lead to generating functions that are not even D-finite, so their interacting counterparts will not be D-finite either. For the rest of the problems, we hope to find the generating functions $Q$ explicitly, or at least the functions $\Qf{0,0}$. This appears to be difficult in most nontrivial cases. Recent research \cite{tabbara2014exact} has shown that there are models, for which we have an explicit relationship between $Q$ or $\Qf{0,0}$ and unmarked variants. As we saw in the previous chapter, the one-dimensional cases admit a linear relationship with polynomial coefficients. There are other two-dimensional models that exhibit a rational relationship. The way things have turned with regards to the \classic models, it would not be surprising if there are two-dimensional interacting models for which neither a linear nor a rational relationship would exist between $Q$, $\Qf{0,0}$ and various unmarked variants. In this section we will cover the first few examples of quarter plane models with interactions. 

Our first step is to find all the reducible models and we can do this  by looking at which of the steps $\{$N, NE, E$\}$ they have in their step set.

\begin{itemize}
	\item[1.] None of the three steps are present. We have $2^5$ such models, in which walks cannot leave the origin and they are reducible to ones with $\mc S = \emptyset$. \awk
	\begin{figure}[H]
		\centering
		\[
			\modelgraphB 10010111
		\]
		\caption{Set of steps allowed in models that do not leave the origin. The models are equivalent to a model without steps.}
	\end{figure}
	\item[2.] There is a $N$ step or an $E$ step. The models are equivalent to one-dimensional ones. There are $2\cdot2^4 = 32$ such models.
	\begin{figure}[H]
		\centering
		\[
			\modelgraphB 11010110 \quad \modelgraphB 00011111
		\]
		\caption{Sets of steps allowed in models with an N step or an E step. The models are essentially one-dimensional.}
	\end{figure}
	\item[3.] There are two or more steps from $\{$N, NE, E$\}$ or NE $\subseteq \mc S$. In this case, the models feature walks that leave both axes and once that happens, such walks may use any step from the step set. Therefore, these models are never reducible.
 \end{itemize}

Cases 1 and 2 contain all of the quarter plane models reducible to one-dimensional models. We will look at the case of a model equivalent to a one-dimensional model along the $x_2$ axis. From the analysis in the previous section, since both $S_1^{\um 1} = 0$ and $S_1^1 = 0$, we have the first situation. If the solution to the one-dimensional problem is $q(x,a;t) = q(x_2,a_2;t)$, then the solution to the models in cases 1 and 2 is $Q(x_1,x_2,a_1,a_2;t) = q(x_2,a_2;a_1t)$. To see it in another way, consider the polynomial $S$ itself in the reduced models. Since $S$ factors, the function $Q$ is a Hadamard product between two one-dimensional models.

Since the models are essentially one-dimensional, we also have a linear relationship using the operators $B_1$ and $B_2$. If $f(a,t)$ is the function that goes together with $q(x,a,t)$, then the function $F(a_1,a_2,t)$ that goes together with $Q(x_1,x_2,a_1,a_2,t)$ is given by
\begin{align}
	F(a_1,a_2,t) = f(a_2,a_1 t)\,.
\end{align}
Then we have $B_2 (F Q) = 0$ and consequently, $B_1B_2(F Q) = 0$. We wrap this section up with the following proposition.
\begin{prop} Each of the 64  reducible models is equivalent to a one-dimensional model confined on the $x_1$ or $x_2$ axis. Their solutions are $Q = q(x_1,a_1;a_2t)$ or $Q = q(x_2,a_2;a_1t)$, respectively, for a generating function $q(x,a;t)$ for a one-dimensional model. Furthermore, they all satisfy a linear relation $B(F Q) = 0$ for operator $B = B_1 B_2$ and function $F$, corresponding to $q$. That is, $F = f(a_2,a_1t)$ or $f(a_1,a_2t)$, respectively. \awk
\end{prop}

%=======================================
\subsection{Kernel Group Considerations}
%=======================================
We can partition the two-dimensional lattice walk models in another way --- with respect to the nature of the group operations in $G_S$. We will combine cases for analysis in the form "case $i.j$": we have case $i$ for $x_1$ and case $j$ for $x_2$, from the section on groups. Of course, case $i.j$ is equivalent to case $j.i$. From the analysis in the previous sections, cases $1.1$, $1.2$ and $1.3$ only feature reducible models, so we may skip them. We are left with the following cases

\begin{itemize}
	\item[Case 2.2:] In this case, both $x_1$ and $x_2$ can only be transformed via the identity transformation. Therefore, exactly one of $S_1^{\um 1}$ and $S_1^1$ is zero and exactly one of $S_2^{\um 1}$ and $S_2^1$ is zero. The models where $S_1^1 = 0$ or $S_2^1 = 0$ already have solutions since they are reducible. The only remaining unsolved case is when $S_1^{\um 1} = 0$ and $S_2^{\um 1} = 0$. The possible models are the ones with steps in $\{$N, E, NE$\}$ except the ones going either only in the N direction or only in the E direction. Here are the models in the positive quadrant which we have not solved yet.
	\begin{figure}[H]
	\centering
	\[
		\modelgraphB 01101000\quad\modelgraphB 01001000\quad \modelgraphB 00101000 \quad \modelgraphB 01100000 \quad \modelgraphB 00100000
	\]
	\caption{Models from case 2.2.}
	\end{figure}
	\item[Case 2.3:] In this case, there is only one transformation for $x_1$ and two for $x_2$. Therefore, $S_1^{\um 1} = 0$ (positive case) or $S_1^{1} = 0$ (negative case) but not both. The negative case is reducible, so we only seek solutions for the positive one. Furthermore, $S_2^{\um 1} \neq 0$ and $S_2^1 \neq 0$, so there must be at least one step going in the positive $x_2$ direction and one in the negative $x_2$ direction. If both NE and SE are in $\mc S$, then we have 8 cases. If one of NE and SE is in $\mc S$ but not both, then we have another 8 cases. Finally, if neither is in $\mc S$, we must have the other three steps present in order to satisfy the conditions. We have 17 cases in total, the group size is two, and none of the models are reducible. Here are all of these models.
	\begin{figure}[H]
		\centering
		\begin{align*}
			\modelgraph 01101011 \quad \modelgraph 00101011 \quad \modelgraph 01100011 \quad \modelgraph 01101001 &\quad \modelgraph 00100011  \quad \modelgraph 00101001  \quad \modelgraph 01100001  \quad \modelgraph 00100001 \\
			\modelgraph 01001011 \quad \modelgraph 01001001 \quad \modelgraph 01000011 \quad \modelgraph 01000001 &\quad \modelgraph 01101010 \quad \modelgraph 01100010 \quad \modelgraph 00101010 \quad \modelgraph 00100010 \quad\modelgraph 01001010
		\end{align*}
		\caption{Models from case 2.3.}
	\end{figure}
	\item[Case 3.3:] This is the most interesting case, with two transformations for $x_1$ and $x_2$ each. Thus, there must be at least one step going in a positive and one going in a negative direction, for each of $x_1$ and $x_2$. If we look at which diagonal steps are present, similarly to case 2.2, we find that there are 145 such models. Previous results show that the group $G(K)$ has orders 4, 6, 8 and infinity, depending on the step set in question. None of these cases are reducible.
 \end{itemize}
 
 %=======================================
\section{Solving Basic Two-dimensional Models}
%=======================================

%=======================================
\subsection{Case 2.2 -- Group of Order One}
%=======================================
The only models that are not reducible are those where $S_1^{\um 1} = 0$ and $S_2^{\um 1} = 0$. Here
\begin{align}
	S = \delta + \alpha x_1 + \beta x_2 + \gamma x_1 x_2
\end{align}
for some $\alpha$, $\beta$, $\gamma$ and $\delta$. The functional equation and its zeroth term extractions are given by
\begin{align}
	\begin{split}
		(1-t S)Q - c_1 Q_1^0- c_2 Q_2^0 + c_1c_2 Q_{1,2}^{0,0}  &= \ol a_1\ol a_2\,,\\
		(1-ta_1 (\delta + \beta x_2))Q_1^0 - c_2Q_{1,2}^{0,0} &=\ol a_2\,,\\
		(1-ta_2 (\delta + \alpha x_1))Q_2^0 - c_1Q_{1,2}^{0,0}  &=\ol a_1\,,\\
		(1-a_1 a_2 t \delta)Q_{1,2}^{0,0} &= 1\,.
	\end{split}
\end{align}
This is a linear system for $\Qf{0,0}$, $Q_1^0$, $Q_2^0$ and $Q$, and by back-substitution,
\begin{align}
	Q &=\frac{1}{1-t S}\left (\ol a_1\ol a_2 + \frac{1}{1-ta_1a_2 \delta}\left(c_1 \frac{1 -  a_1 \delta t}{1-t a_1S_1^0} + c_2 \frac{1 -  a_2 \delta t}{1-t a_2S_2^0} - c_1 c_2 \right) \right)\\
	&= \frac{1}{1-t S}\left(\ol a_1\ol a_2 + \frac{1}{1-a_1a_2\delta t}(c_1 f_1(a_1,t) + c_2 f_2(a_2,t)-c_1c_2)\right)\,.
\end{align}
Note how $Q$ depends on functions $f_i$ that are each missing one of the $a$ variables. We shall see that the linear relation $B(F Q) = 0$ is satisfied by the function
\begin{gather}
	\begin{align}
		F(a_1,a_2,t) = 1 - a_1 a_2 \delta t\,.
	\end{align}
	\shortintertext{First,}
	\begin{align}
		B((1 - a_1 a_2 \delta t)\ol a_1\ol a_2) &= \ol a_1\ol a_2- \ol a_2-\ol a_1 + 1 = c_1 c_2\,,\\
		B(c_1 f_1 + c_2 f_2-c_1c_2) &= c_1 f_1 + c_2 f_2 -c_1c_2 - c_2 f_2 - c_1 f_1 = -c_2 c_2\,.
	\end{align}
\end{gather}
Putting the two quantities together, along with dividing by $1-t S$, shows that $B(F Q) = 0$. Of course, this works for $\Qf{0,0}$ too.
%=======================================
\subsection{Case 2.3 -- Group of Order Two}
%=======================================
In this case only $S_1^{\um 1}$ is zero, so the functional equation are
\begin{align}
	\begin{split}
		(1-t S)Q - c_1 Q_1^0- (c_2 -t\ol x_2 S_1^{\um 1})Q_2^0 + c_1c_2 Q_{1,2}^{0,0}  &= \ol a_1\ol a_2\,,\\
		(1-ta_1 S_1^0)Q_1^0 - (c_2 - ta_1 \ol x_2 \Sf{0,\um 1})Q_{1,2}^{0,0} &=\ol a_2\,,\\
		(1-ta_2 S_2^0)Q_2^0 - c_1Q_{1,2}^{0,0}  - a_2 t S_2^{\um 1} Q_2^1&=\ol a_1\,,\\
		(1-a_1 a_2 t \Sf{0,0})Q_{1,2}^{0,0} - ta_1 a_2 \Sf{0,\um 1} \Qf{0,1} &= 1\,.
	\end{split}
\end{align}
We can identify a subproblem that is essentially one-dimensional and therefore, already solved. To see this, take the second and fourth equations, and make the six substitutions $a_i = a$, $c_i = c$, $x_2 = x$, $a_1 t = \tau$, $S_1^0 = T$ and $Q_1^0 (x_2)= q(x)$. The two equations become
\begin{align}
	(1-\tau T)q - (c - \tau  \ol x T_{\um 1})q_0 &=\ol a\,,\\
	(1-a \tau T_0)q_0 - a\tau T_{\um 1} q_1 &= 1\,.
\end{align}
We can see that this is a one-dimensional problem again. The solution here is $q(x,a;\tau)$, obtained by the methods in the previous chapter, so the original problem has solutions
\begin{align}
	Q_1^0(x_2) = q(x_2,a_2;a_1t)\quad \text{and}\quad Q_{1,2}^{0,0} = q_0(a_2;a_1t)\,.
\end{align}
Since $\Qf{0,0}$ and $Q_1^0$ are given by a one-dimensional model's solution, they have a linear relationship with their partially-marked versions. Next, we rewrite the master equation as
\begin{align}
	(1-t S)Q - (c_2 -t\ol x_2 S_2^{\um 1})Q_2^0  &= \ol a_1(\ol a_2 + b_1 q  -b_1c_2q_0)  =\ol a_1\phi(x_2)\,.
\end{align}
Since $S_1^{\um 1} \neq 0$ and $S_2^1 \neq 0$, we have a root $x_2 = \sigma_2$ that allows us to eliminate the $(1-t S)Q$ term. That is, we are finding a root for $x_2$ again, since $x = x_2$ in the previous equations. We find that
\begin{gather}
	\begin{align}
		Q_2^0  &=-\frac{\ol a_1\phi(\sigma_2)}{ c_2 -t\ol \sigma_2 S_2^{\um 1}}
	\end{align}
	\shortintertext{and}
	\begin{align}
		Q = \frac{1}{1-t S}\left(\ol a_1\phi(x_2) - \frac{c_2 -t\ol x_2 S_2^{\um 1}}{c_2 -t\ol \sigma_2 S_2^{\um 1}}\ol a_1\phi(\sigma_2)\right)\,.
	\end{align}
\end{gather}
We may simplify this form based on what $q$ and $q_0$ are from the one-dimensional setting. We will do something else. After a quick computation, we have $Q_2^0 = \ol a_1 a_2 \phi(\sigma_2) \wh Q_2^0$ and we may write
\begin{align}
	KQ &= \ol a_1 \phi + K_2\ol a_1 a_2 \phi(\sigma_2)\wh Q_2^0\,,\\
	K\wh Q &= \ol a_2 + K_2\wh Q_2^0\,.	
\end{align}
Therefore, after expanding the definitions of the various quantities, we have
\begin{align}
	KQ - \ol a_1 a_2 \phi(\sigma_2)K\wh Q &= \ol a_1 (\phi - \phi(\sigma_2)) = c_1 (q-q(\sigma_2))
\end{align}
Since $q = q(x_2,a_2;a_1t)$, there is a function $f$ depending on $a$'s and $t$ only, for which $B_2(fq) = 0$. The same function works for $q(\sigma_2)$ as well. Finally, there is a linear relationship between $Q$ and its various partially-interacting versions. This time it seems that the relationship is not given by $B(F Q) = 0$ but rather a more complicated one, $B_2(fg(B_1(Q/g))) = 0$, where $g = \ol a_1 a_2 \phi(\sigma_2)$. Finally, we may give the following theorem.

\begin{thm} Consider the two-dimensional \classic models. If the model is reducible or has a group of order one or two, then the generating function depends on its partially-interacting versions in a linear manner.
\end{thm}

%=======================================
\subsection{Miscellaneous Case}
%=======================================

There is one case with an infinite group that is both easy to solve and not reducible --- the model with $S = x_1x_2 + \ol x_1 \ol x_2$. The transformations are $g_1= \ol x_1 \ol x_2^2$ and $g_2= \ol x_2 \ol x_1^2$. If we apply them on $(x_1,x_2)$ in an alternating sequence, each time the absolute powers of $x_1$ and $x_2$ will increase in one of the coordinates, and therefore, the group is infinite. However, this case is easy to solve because the model is essentially one-dimensional. It is equivalent to the model on the non-negative line with $S = x + \ol x$ and marking visits to the origin with $a$, where $x = x_1 x_2$ and $a = a_1 a_2$. The solution is
\begin{align}
	Q(x_1,x_2,a_1,a_2;t) = q(x_1x_2,a_1a_2;t)
\end{align}
and we can add this model to the theorem as well.

%=======================================
\subsection{Counting Classes of Models}
%=======================================

We will tally all the models on the quarter plane. We have 64 reducible models, 5 new cases with group of order one and 34 new cases with group of order two. Also, there are 29 models with group of order four, 6 with order six, and 4 with order eight, plus one easy case with infinite group. So far we have 143 cases of models we can solve easily or models with a finite group, some of which we can solve easily. Next, by inspection of Table 4 in \cite[p.~28]{bousquet2010walks}, we see there are 120 models with infinite groups: 12 with three steps, 40 with four steps, 44 with five steps, and 24 with six steps. Adding all of those together yields 263 cases. There are 7 reducible models with an infinite group that we have counted twice. The total count is 256 and we are done.
\begin{figure}[H]
	\centering
	\begin{align*}
		\modelgraph 00100100 \quad \modelgraph 10000011 \quad \modelgraph 10010001 \quad \modelgraph 10000101 \quad \modelgraph 10010011 \quad \modelgraph 10000111 \quad \modelgraph 10010101 \quad \modelgraph 10010111
	\end{align*}
	\caption{Reducible models with an infinite group.}
\end{figure}

So far we have solved the 64 reducible models, the 39 other models with group of order one or two, and the single new case of easy-to-solve model with an infinite group, for a grand total of 104 trivial or easy models. Out of the remaining 152 models, 113 have an infinite group, and their analysis will not be a part of this thesis. Finally, we have 39 cases left, with a group of order four, six or eight. The rest of the thesis will handle the analysis of those cases.

\begin{figure}[H]
	\centering
	\begin{align*}
		\modelgraph 01011010 \quad \modelgraph 10100101  \quad \modelgraph 11100111 \quad \modelgraph 11111111 \quad \modelgraph 10100010 \quad &\modelgraph 10111010 \quad \modelgraph 11111010 \quad \modelgraph 11100010  \quad \modelgraph 11100101  \quad \modelgraph 11111101\\
	 	\modelgraph 00110010 \quad \modelgraph 11011011 \!\exxxx&\, \quad\quad \modelgraph 10011001
	\end{align*}
	\caption{Models with a finite group.}
\end{figure}

\begin{figure}[H]
	\centering
	\begin{align*}
		\modelgraph 11000001 \quad \modelgraph 10100001 \quad\modelgraph 10011010  \quad \modelgraph 00110011 &\quad \modelgraph 00110110 \quad \modelgraph 10110001 \quad \modelgraph 11001001 \quad \modelgraph 11100001\\
		\modelgraph 11110001 \quad \modelgraph 11101010 \quad \modelgraph 11101001 \quad\, \quad&\!\exxxx \modelgraph 01111010 \quad \modelgraph 11001011 \quad \modelgraph 10111100 \quad \modelgraph 10110011 \\
		\modelgraph 11111001 \quad \modelgraph 11110011 &\quad \modelgraph 10110111 \quad \modelgraph 11111011
	\end{align*}
	\caption{Models with an infinite group. To get the rest of them, apply one of the symmetries of the square.}
\end{figure}

%=====================================================================================
%=====================================================================================
\endinput

%\begin{figure}[H]
%    \centering
%    \includegraphics[scale = 0.5]{placeholder}
%    \caption{Two-dimensional lattice walk model constraints.}
%    \label{fig:plane}
%\end{figure}
%\begin{figure}[H]
%    \centering
%    \includegraphics[scale = 0.5]{placeholder}
%    \caption{Inclusion-exclusion of regions.}
%    \label{fig:inclusion}
%\end{figure}
